\subsection{Normal endomorphisms}

DEFINITION 4. --- Let E be a hilbertian space and \(u \in \mathcal{L}(\mathrm{E})\) . We say that \(u\) is normal if it commutes with its adjoint \(u^{*}\) .

For example, every automorphism \(u\) of the hilbertian space \(E\) is normal since we have \(uu^{*} = u^{*}u = 1_{E}\) .

PROPOSITION 7. --- For \(u \in \mathcal{L}(\mathrm{E})\) to be normal, it is necessary and sufficient that \(\| u(x) \| = \| u^{*}(x) \|\) for all \(x \in \mathrm{E}\) .

We define a hermitian form \(\Phi\) on \(E\) by

\[
\Phi (x, y) = \langle u u ^ {*} (x) | y \rangle - \langle u ^ {*} u (x) | y \rangle .
\]

For u to be normal, it is necessary and sufficient that \(\Phi = 0\) . By the polarization formulas (V, p.~2), this is equivalent to \(\Phi(x, x) = 0\) for all \(x \in E\) . The proposition now follows since

\[
\Phi (x, x) = \left\| u ^ {*} (x) \right\| ^ {2} - \left\| u (x) \right\| ^ {2}.
\]

PROPOSITION 8. --- Suppose that \(u \in \mathcal{L}(\mathrm{E})\) is normal. Let \(\mathbf{N}\) be the kernel of \(u\) and \(\mathbf{M}\) the orthogonal of \(\mathbf{N}\) in \(\mathrm{E}\) ; let \(m\) and \(n\) be two positive integers such that \(m + n \geqslant 1\) . Then \(\mathbf{N}\) is the kernel of \(u^{m}(u^{*})^{n}\) and \(\mathbf{M}\) is both the initial and the final subspace of \(u^{m}(u^{*})^{n}\) . In particular, \(\mathbf{M}\) is both the initial and the final subspace of \(u\) and of \(u^{*}\) , and is stable under \(u\) and \(u^{*}\) .

Prop. 7 shows that u and \(u^{*}\) have the same kernel N. By prop. 4, (ii) of V, p.~41, the subspace M of E is stable under u and \(u^{*}\) since this is so for \(N = M^{\circ}\) , since \(M \cap N = \{0\}\) , the endomorphisms of M induced by u and \(u^{*}\) are injective. Let \(v = u^{m}(u^{*})^{n}\) ; the preceding argument shows that the restriction of v to M (resp. N) is injective (resp. null), hence N is the kernel of v. Consequently, \(M = N^{\circ}\) is the initial subspace of v. By prop. 4, (i) of V, p.~41, the final subspace of v is equal to the initial subspace of \(v^{*}\) . But \(v^{*} = u^{n}(u^{*})^{m}\) and so the initial subspace of \(v^{*}\) is equal to M by the preceding.

COROLLARY. --- Let \(\lambda \in \mathbf{K}\) . The following subspaces of E are equal :

\begin{enumerate}
\def\labelenumi{\alph{enumi})}
\item
  the eigen subspace of \(u\) relative to \(\lambda\) ;
\item
  the eigen subspace of \(u^{*}\) relative to \(\overline{\lambda}\) ;
\item
  the primary subspace of \(u\) relative to \(\lambda\) (in other words, by LIE, VII, § 1, No.~1, the set of all vectors \(x\) of \(E\) for which there exists an integer \(n \geqslant 0\) such that \((u - \lambda .1_{E})^{n}(x) = 0\) );
\item
  the primary subspace of \(u^{*}\) relative to \(\overline{\lambda}\) .
\end{enumerate}

It is clear that \(w = u - \lambda \cdot 1_{E}\) is a normal endomorphism of E, hence the endomorphisms w, \(w^{*} = u^{*} - \overline{\lambda} \cdot 1_{E}\) , \(w^{n}\) and \((w^{*})^{n}\) of E have the same kernel by prop. 8.

\subsection{Hermitian endomorphisms}

DEFINITION 5. --- Let E be a hilbertian space and let \(u \in \mathcal{L}(\mathrm{E})\) . We say that \(u\) is hermitian if \(u^* = u\) .

Let \(\mathcal{H}(\mathrm{E})\) denote the set of all hermitian elements of \(\mathcal{L}(\mathrm{E})\) ; this is a vector subspace of the vector space \(\mathcal{L}(\mathrm{E})_{[\mathbb{R}]}\) over \(\mathbf{R}\) which is deduced from \(\mathcal{L}(\mathrm{E})\) by restricting the scalars.

To each \(u \in \mathcal{L}(\mathrm{E})\) , we associated (V, p.~16, cor. 2) a sesquilinear form \(\Phi_u : (x, y) \mapsto \langle x | u(y) \rangle\) on \(\mathbf{E} \times \mathbf{E}\) . We have

\[
\Phi_ {u ^ {*}} (x, y) = \overline {{{{\Phi_ {u} (y , x)}}}} \quad (x, y \text {in E});\tag{14}
\]

consequently, \(u\) is hermitian if and only if the form \(\Phi_u\) is hermitian. When \(K\) is \(C\) , it is enough to assume that \(\Phi_u(x, x) = \langle x | u(x) \rangle\) is real for all \(x \in E\) (V, p.~2, Remark).

Let \(u \in \mathcal{L}(\mathrm{E})\) . We have seen (V, p.~16, cor. 2) that the norm of \(u\) can be calculated by the formula

\[
\| u \| = \sup _ {\| x \| \leqslant 1, \| y \| \leqslant 1} \left| \Phi_ {u} (x, y) \right|.\tag{15}
\]

When u is hermitian, we have the following result :

PROPOSITION 9. --- For every hermitian endomorphism u of E, we have

\[
\| u \| = \sup _ {\| x \| \leqslant 1} | \langle x | u (x) \rangle |.\tag{16}
\]

Put \(\Phi = \Phi_{u}\) and \(c = \sup_{\| x\| \leqslant 1}|\Phi (x,x)|\) , then evidently \(c\leqslant \| u\|\) . Let \(x,y\) be in E such that \(\| x\| \leqslant 1\) , \(\| y\| \leqslant 1\) . Then

\[
\Phi (x + y, x + y) = \Phi (x, x) + \Phi (y, y) + 2 \mathscr {R} \Phi (x, y),
\]

hence

\[
4 \mathscr {R} \Phi (x, y) = \Phi (x + y, x + y) - \Phi (x - y, \dot {x} - y);
\]

but \(|\Phi(t, t)| \leqslant c \|t\|^2\) for all \(t \in \mathbf{E}\) , thus

\[
4 \left| \mathcal {R} \Phi (x, y) \right| \leqslant c \left(\| x + y \| ^ {2} + \| x - y \| ^ {2}\right) = 2 c \left(\| x \| ^ {2} + \| y \| ^ {2}\right) \leqslant 4 c.
\]

Let \(a = \Phi(x, y)\) ; there exists a complex number \(\lambda\) with absolute value 1 such that \(\lambda a = |a|\) . Replacing \(y\) by \(\lambda y\) in the preceding inequality, we get \(|\Phi(x,y)| \leqslant c\) . By (15), \(\|u\| \leqslant c\) and the proposition follows. Q.E.D.

Evidently every hermitian endomorphism is normal. Conversely :

PROPOSITION 10. --- Suppose K is C. Let \(u \in \mathcal{L}(\mathrm{E})\) . Then there exists a unique pair \((h_1, h_2)\) of hermitian endomorphisms of E, such that \(u = h_1 + ih_2\) . In order that u is normal, it is necessary and sufficient that \(h_1\) and \(h_2\) commute.

For, the relation \(\ll u = h_1 + ih_2\) , \(h_1^* = h_1\) , \(h_2^* = h_2\) is equivalent to

\[
\ll h _ {1} = \frac {1}{2} (u + u ^ {*}) \text { and } h _ {2} = \frac {i}{2} (u ^ {*} - u) \gg .
\]

In addition, we have \(h_1h_2 - h_2h_1 = \frac{i}{2} (uu^* - u^*u)\) . This proves prop. 10.

PROPOSITION 11. --- Let \(p \in \mathcal{L}(\mathbf{E})\) . In order that \(p\) is the orthoprojector from \(\mathbf{E}\) onto a closed vector subspace of \(\mathbf{E}\) , it is necessary and sufficient that \(p^2 = p = p^*\) .

Suppose \(p^{2} = p\) . Let M be the image of p and N its kernel. E is the topological direct sum of M and N. In order that p is an orthoprojector, it is necessary and sufficient that M is orthogonal to N, that is to say that we have \(\langle p(x)|y - p(y)\rangle = 0\) for all x, y in E. This latter relation is equivalent to \(p = p^{*}p\) , and implies that \(p^{*} = (p^{*}p)^{*} = p^{*}p = p\) ; conversely if \(p^{*} = p\) , we have \(p = p^{2} = p^{*}p\) .

\subsection{Positive endomorphisms}

DEFINITION 6. --- Let E be a hilbertian space and \(u \in \mathcal{L}(E)\) . We say that u is positive, and write \(u \geqslant 0\) , if u is hermitian and if \(\langle x|u(x) \rangle \geqslant 0\) for all \(x \in E\) .

When K is equal to C, the relation

\[
\langle x | u (x) \rangle \geqslant 0 \quad \text { for   all } \quad x \in E
\]

implies that \(u\) is hermitian (V, p.~2, Remark), hence positive.

Let \(\mathcal{L}_{+}(\mathrm{E})\) denote the set of all positive elements of \(\mathcal{L}(\mathrm{E})\) ; this is a proper pointed convex cone in the real vector space \(\mathcal{L}(\mathrm{E})_{[\mathbb{R}]}\) underlying \(\mathcal{L}(\mathrm{E})\) . In order that \(u\) is positive, it is necessary and sufficient that the sesquilinear form \(\Phi_u\) on \(\mathrm{E} \times \mathrm{E}\) associated with \(u\) is positive hermitian. Given \(u\) and \(v\) in \(\mathcal{L}(\mathrm{E})\) , the relation \(u - v \geqslant 0\) can also be written as \(u \geqslant v\) or \(v \leqslant u\) ; this is an order relation on \(\mathcal{L}(\mathrm{E})_{[\mathbb{R}]}\) compatible with its real vector space structure.

PROPOSITION 12. --- Let u be a hermitian (resp. positive) element of \(\mathcal{L}(\mathrm{E})\) and let v be a continuous linear mapping from E into a hilbertian space F. Then vuv* is a hermitian (resp. positive) element of \(\mathcal{L}(\mathrm{F})\) .

For, we have \((vuv^{*})^{*} = v^{**}u^{*}v^{*} = vuv^{*}\) . On the other hand, if \(u \geqslant 0\) , we have

\[
\langle y | v u v ^ {*} (y) \rangle = \langle v ^ {*} (y) | u (v ^ {*} (y)) \rangle \geqslant 0
\]

for all \(y \in \mathbf{F}\) , hence \(vuv^{*} \geqslant 0\) .

Prop. 12 shows, in particular, that \(vv^*\) is positive for all \(v \in \mathcal{L}(\mathrm{E};\mathrm{F})\) . Since, in particular, an orthoprojector \(p\) satisfies \(p = p^2 = pp^*\) , it is positive.

Remarks. --- 1) For every hermitian \(u\) in \(\mathcal{L}(\mathrm{E})\) , put \(m(u) = \inf_{||x|| = 1} \langle x|u(x) \rangle\) , \(\mathbf{M}(u) = \sup_{||x|| = 1} \langle x|u(x) \rangle\) . If E is not just 0, \(m(u)\) and \(\mathbf{M}(u)\) are finite; moreover, \(\mathbf{M}(u)\) is the smallest real number \(\lambda\) such that \(u \leqslant \lambda\) . \(1_{\mathrm{E}}\) and \(m(u)\) the largest real number \(\mu\) such that \(u \geqslant \mu\) . \(1_{\mathrm{E}}\) . Clearly we have \(m(-u) = -\mathbf{M}(u)\) and \(\mathbf{M}(-u) = -m(u)\) . It is clear that

\[
\sup \left(| m (u) |, | \mathrm{M} (u) |\right) = \sup _ {\| x \| = 1} \left| \langle x | u (x) \rangle \right|
\]

and prop. 9 (V, p.~44) implies (for E ≠ \{0\}) that

\[
\| u \| = \sup \left(| m (u) |, | \mathbf {M} (u) |\right).\tag{17}
\]

* For another proof of this formula when K is C, see prop. 14 of TS, I, § 6, No.~8. * 2) Let M and N be two closed vector subspaces of E, and \(p_{\mathbf{M}}\) (resp. \(p_{\mathbf{N}}\) ) the orthoprojector from E onto M (resp. N). Then \(\mathbf{M} \subset \mathbf{N}\) if and only if \(p_{\mathbf{M}} \leqslant p_{\mathbf{N}}\) . For, we have \(p_{\mathbf{M}}^{*}p_{\mathbf{M}} = p_{\mathbf{M}}\) , hence

\[
\left\| p _ {\mathbf {M}} (x) \right\| ^ {2} = \langle p _ {\mathbf {M}} (x) | p _ {\mathbf {M}} (x) \rangle = \langle x | p _ {\mathbf {M}} ^ {*} p _ {\mathbf {M}} (x) \rangle = \langle x | p _ {\mathbf {M}} (x) \rangle
\]

for all \(x \in E\) . The relation \(p_{M} \leqslant p_{N}\) is therefore equivalent to « \(\|p_{M}(x)\| \leqslant \|p_{N}(x)\|\) for all \(x \in E\) ». If \(M \subset N\) , we have \(p_{M} = p_{M} p_{N}\) , hence \(\|p_{M}(x)\| \leqslant \|p_{N}(x)\|\) since \(\|p_{M}\| \leqslant 1\) . Conversely, if \(\|p_{M}(x)\| \leqslant \|p_{N}(x)\|\) for all \(x \in E\) , the kernel of \(p_{M}\) contains the kernel of \(p_{N}\) , that is, that \(M^{\circ} \supset N^{\circ}\) , which implies that \(M \subset N\) .

PROPOSITION 13. --- Let \(\mathcal{H}(\mathrm{E})\) be the set of all continuous hermitian endomorphisms of the hilbertian space E. Let \(\mathcal{F}\) be a non-empty, directed increasing and bounded subset of \(\mathcal{H}(\mathrm{E})\) .

\begin{enumerate}
\def\labelenumi{(\roman{enumi})}
\tightlist
\item
  The set \(\mathcal{F}\) has an upper bound \(u_0\) in \(\mathcal{H}(\mathrm{E})\) ; we have
\end{enumerate}

\[
\langle x | u _ {0} (x) \rangle = \sup _ {u \in \mathcal {F}} \langle x | u (x) \rangle \quad f o r a l l \quad x \in E.\tag{18}
\]

\begin{enumerate}
\def\labelenumi{(\roman{enumi})}
\setcounter{enumi}{1}
\tightlist
\item
  The filter of sections of \(\mathcal{F}\) converges to \(u_0\) in the space \(\mathcal{L}(\mathrm{E})\) endowed with the topology of simple convergence.
\end{enumerate}

Let \(\Sigma\) be the filter of sections of \(\mathcal{F}\) ; for every \(u \in \mathcal{H}(\mathrm{E})\) , let \(\Phi_u\) be the continuous hermitian form on \(\mathbf{E}\) defined by

\[
\Phi_ {u} (x, y) = \langle x | u (y) \rangle .
\]

Let

\[
\Psi_ {u} (x) = \Phi_ {u} (x, x)
\]

for \(u \in \mathcal{H}(\mathrm{E})\) and \(x \in \mathrm{E}\) . By the polarization formulas (V, p.~2), we have

\[
4 \Phi_ {u} (x, y) = \Psi_ {u} (x + y) - \Psi_ {u} (x - y) \quad \text { if } \quad \mathbf {K} = \mathbf {R}\tag{19}
\]

\[
4 \Phi_ {u} (x, y) = \Psi_ {u} (x + y) - \Psi_ {u} (x - y) - i \Psi_ {u} (x + i y) + i \Psi_ {u} (x - i y) \quad \text { if } \quad \mathbf {K} = \mathbf {C}.\tag{20}
\]

For every \(x \in \mathbf{E}\) , the mapping \(u \mapsto \Psi_u(x)\) from into \(\mathbf{R}\) is increasing and bounded, hence has a limit with respect to \(\Sigma\) . By the preceding formulas, the limit

\[
\lim _ {u, \Sigma} \Phi_ {u} (x, y) = \Phi (x, y)
\]

exists for every pair \((x, y)\) of elements of E. It is clear that \(\Phi\) is a hermitian form on E. If \(v_{1} \in F\) and \(v_{2}\) is a bound of F, the hermitian forms \(f_{1} = \Phi - \Phi_{v_{1}}\) and \(f_{2} = \Phi_{v_{2}} - \Phi\) are positive; there exists a real number \(M \geqslant 0\) such that

\[
f _ {1} (x, x) + f _ {2} (x, x) = \Phi_ {v _ {2} - v _ {1}} (x, x) \leqslant \mathbf {M} \| x \| ^ {2},
\]

hence

\[
f _ {1} (x, x) \leqslant \mathbf {M} \| x \| ^ {2}, f _ {2} (x, x) \leqslant \mathbf {M} \| x \| ^ {2} (x \in \mathrm{E});
\]

consequently the semi-norms \(x \mapsto f_i(x, x)^{1/2}\) are continuous on E. Since

\[
f _ {2} - f _ {1} = \Phi_ {v _ {2}} + \Phi_ {v _ {1}} - 2 \Phi ,
\]

we conclude that \(x \mapsto \Phi(x, x)\) is a continuous function on E, and by formulas (19) and (20), that \(\Phi\) is continuous on \(E \times E\) . Therefore there exists (V, p.~16, cor. 2) an element \(u_{0}\) of \(\mathcal{H}(E)\) such that \(\Phi = \Phi_{u_{0}}\) . Formula (18) is evidently satisfied, hence \(u_{0}\) is the upper bound of F in \(\mathcal{H}(E)\) . This proves (i).

We have, by construction

\[
\lim _ {u, \Sigma} \langle x | (u _ {0} - u) (x) \rangle = 0 \quad \text { for   all } \quad x \in E.\tag{21}
\]

Let \(v_{1} \in \mathcal{F}\) ; given a \(u \in \mathcal{F}\) such that \(u \geqslant v_{1}\) , let \(v = u_{0} - u\) . If we apply the Cauchy-Schwarz inequality to the positive hermitian form \(\Phi_v\) on E, we get

\[
\begin{array}{r l} \left\| v (x) \right\| ^ {4} & = \left| \Phi_ {v} (v (x), x) \right| ^ {2} \leqslant \Phi_ {v} (v (x), v (x)). \Phi_ {v} (x, x) \\ & = \langle v (x) | v ^ {2} (x) \rangle \langle x | v (x) \rangle \leqslant \| v \| ^ {3} \| x \| ^ {2} \langle x | v (x) \rangle \\ & \leqslant \| u _ {0} - v _ {1} \| ^ {3} \| x \| ^ {2} \langle x | v (x) \rangle , \end{array}
\]

since \(\| v\| \leqslant \| u_0 - v_1\|\) by V, p.~44, prop. 9. Then by (21) we get \(\lim_{u,\Sigma}\left\| (u_0 - u)(x)\right\| = 0\) for all \(x\in \mathrm{E}\) ; which proves assertion (ii). Q.E.D.

In particular, prop. 13 can be applied to the case of an increasing and bounded sequence \((u_n)_{n\in \mathbf{N}}\) of elements of \(\mathcal{H}(\mathrm{E})\) . Then there exists an element \(v\) of \(\mathcal{H}(\mathrm{E})\) characterized by

\[
\langle x | v (x) \rangle = \lim _ {n \rightarrow \infty} \langle x | u _ {n} (x) \rangle = \sup _ {n \in \mathbf {N}} \langle x | u _ {n} (x) \rangle \quad (x \in E),
\]

and we have \(v(x) = \lim_{n\to \infty}u_n(x)\) for all \(x\in E\) . Moreover, \(v\) is the upper bound of the set of the \(u_{n}\) in \(\mathcal{H}(E)\) .

\subsection{Trace of an endomorphism}

Let E and F be two hilbertian spaces. Conforming to the conventions of V, p.~40, we let \(ba^{*}\) , for \(a\) in E and \(b\) in F, denote the continuous linear mapping \(x \mapsto b \langle a|x\rangle\) from E into F.

Lemma 1. --- There exists an isomorphism \(\theta\) from the vector space \(\mathbf{F} \otimes \mathbf{E}'\) onto the space \(\mathcal{L}_f(\mathbf{E}; \mathbf{F})\) of all finite rank continuous linear mappings from \(\mathbf{E}\) into \(\mathbf{F}\) , characterized by \(\theta(b \otimes a^*) = ba^*\) for \(a \in \mathbf{E}\) , \(b \in \mathbf{F}\) .

By A, II, § 4, No.~2, there exists an injective linear mapping \(\theta\) from \(\mathbf{F} \otimes \mathbf{E}'\) into \(\mathcal{L}(\mathbf{E};\mathbf{F})\) and only one such, which transforms \(b \otimes a'\) into the linear mapping \(x \mapsto ba'(x)\) for \(a' \in \mathbf{E}'\) , \(b \in \mathbf{F}\) . Evidently \(\theta(b \otimes a^*) = ba^*\) , and the image of \(\theta\) is contained in \(\mathcal{L}_f(\mathbf{E};\mathbf{F})\) . However, let \(u \in \mathcal{L}_f(\mathbf{E};\mathbf{F})\) and let \((e_1,\dots,e_n)\) be an orthonormal basis of the image of \(u\) in \(\mathbf{F}\) . Let \(f_i = u^*(e_i)\) for \(1 \leqslant i \leqslant n\) . For every \(x \in \mathbf{E}\) , we have

\[
u (x) = \sum_ {i = 1} ^ {n} \left\langle e _ {i} | u (x) \right\rangle . e _ {i} = \sum_ {i = 1} ^ {n} \left\langle f _ {i} | x \right\rangle . e _ {i},
\]

hence \(u = \sum_{i=1}^{n} e_i f_i^* = \theta(\sum_{i=1}^{n} e_i \otimes f_i^*)\) . Therefore the image of \(\theta\) is equal to \(\mathcal{L}_f(\mathrm{E};\mathrm{F})\) .

\[
\mathrm{E} = \mathrm{F},
\]

\[
\mathscr {L} _ {f} (\mathrm{E}) = \mathscr {L} _ {f} (\mathrm{E}; \mathrm{E})
\]

\[
\mathcal {L} _ {f} (\mathrm{E})
\]

\[
\tau (\theta (a \otimes a ^ {\prime})) = a ^ {\prime} (a)
\]

\[
a \in \mathrm{E}, a ^ {\prime} \in \mathrm{E} ^ {\prime}\tag{22}
\]

\[
\tau (b a ^ {*}) = \langle a | b \rangle \quad \text { for } a, b \text { in } E.
\]

When E is finite dimensional, we have \(\mathcal{L}_{f}(\mathrm{E}) = \mathcal{L}(\mathrm{E})\) and \(\tau(u)\) is the trace of the endomorphism u of E (A, II, § 4, No.~3).

Lemma 2. --- Let \((e_i)_{i\in \mathbf{I}}\) be an orthonormal basis of E. Then

\[
\tau (u) = \sum_ {i \in \mathrm{I}} \left\langle e _ {i} | u (e _ {i}) \right\rangle
\]

for all \(u \in \mathcal{L}_f(\mathrm{E})\) .

It is enough to consider the case when \(u = ba^{*}\) with \(a, b\) in E. Then

\[
\langle e _ {i} | u (e _ {i}) \rangle = e _ {i} ^ {*} b. a ^ {*} e _ {i} = \overline {{\langle e _ {i} | a \rangle}} \langle e _ {i} | b \rangle
\]

and lemma 2 follows from formula (22) and formula (3) of V, p.~22.

Lemma 3. --- Let u be a continuous and positive endomorphism of E, and F the set of all finite rank orthoprojectors on E. Then for every orthonormal basis \((e_{i})_{i\in\mathbf{I}}\) of E, we have \((in\overline{\mathbf{R}}_{+})\) the equality

\[
\sum_ {i \in \mathrm{I}} \left\langle e _ {i} | u (e _ {i}) \right\rangle = \sup _ {p \in \mathcal {F}} \tau (p u p).
\]

For every finite subset J of I, put \(p_{J} = \sum_{i \in J} e_{i} e_{i}^{*}\) ; this is the orthoprojector from E onto the vector subspace generated by the vectors \(e_{i}\) , where i ranges over J. We have

\[
p _ {\mathrm{J}} u p _ {\mathrm{J}} = \sum_ {i \in \mathrm{J}, j \in \mathrm{J}} \left\langle e _ {i} | u (e _ {j}) \right\rangle e _ {i} e _ {j} ^ {*},
\]

hence \(\tau(p_{\mathbf{J}}up_{\mathbf{J}}) = \sum_{i\in \mathbf{J}}\langle e_i|u(e_i)\rangle\) . Since \(p_{\mathbf{J}}\in \mathcal{F}\) ,

\[
\sum_ {i \in \mathbf {J}} \left\langle e _ {i} | u (e _ {i}) \right\rangle \leqslant \sup _ {p \in \mathcal {F}} \tau (p u p);
\]

and so we conclude that

\[
\sum_ {i \in \mathbf {I}} \left\langle e _ {i} | u (e _ {i}) \right\rangle = \sup _ {\mathbf {J}} \sum_ {i \in \mathbf {J}} \left\langle e _ {i} | u (e _ {i}) \right\rangle \leqslant \sup _ {p \in \mathcal {F}} \tau (p u p).
\]

Let \(v\) be a finite rank continuous and positive endomorphism of \(\mathbf{E}\) and let \(p \in \mathcal{F}\) . By th. 2 of \(\mathbf{V}\) , p.~23 there exists an orthonormal basis \((f_{\alpha})_{\alpha \in \mathbf{A}}\) of \(\mathbf{E}\) and a finite subset \(\mathbf{B}\) of \(\mathbf{A}\) such that \((f_{\alpha})_{\alpha \in \mathbf{B}}\) is an orthonormal basis of the image of \(p\) . Then we have \(p = \sum_{\alpha \in \mathbf{B}} f_{\alpha} f_{\alpha}^{*}\) , and so, as above, the relation \(\tau(pvp) = \sum_{\alpha \in \mathbf{B}} \langle f_{\alpha}|v(f_{\alpha}) \rangle\) . By lemma 2 (V, p.~48) we have \(\tau(v) = \sum_{\alpha \in \mathbf{A}} \langle f_{\alpha}|v(f_{\alpha}) \rangle\) , which gives the formula

\[
\sum_ {\alpha \in \mathbf {B}} \langle f _ {\alpha} | v (f _ {\alpha}) \rangle \leqslant \tau (v).
\]

Applying this inequality to the case where \(v = p_{\mathbf{J}}.up_{\mathbf{J}}\) and where \(\mathbf{J}\) is a finite subset of \(\mathbf{I}\) , we get

\[
\sum_ {\alpha \in \mathbf {B}} \left\langle p _ {\mathbf {J}} (f _ {\alpha}) | u p _ {\mathbf {J}} (f _ {\alpha}) \right\rangle \leqslant \sum_ {i \in \mathbf {J}} \left\langle e _ {i} | u (e _ {i}) \right\rangle .\tag{23}
\]

For every \(x \in \mathbb{E}\) , we have \(p_{\mathbf{J}}(x) = \sum_{i \in \mathbf{J}} \langle e_i | x \rangle e_i\) , and so \(x = \lim_{\mathbf{J}} p_{\mathbf{J}}(x)\) with respect to the ordered directed set of finite subsets \(\mathbf{J}\) of \(\mathbf{I}\) . Passing to the limit over \(\mathbf{J}\) in (23), we get

\[
\tau (p u p) = \sum_ {\alpha \in \mathrm{B}} \left\langle f _ {\alpha} | u (f _ {\alpha}) \right\rangle \leqslant \sum_ {i \in \mathrm{I}} \left\langle e _ {i} | u (e _ {i}) \right\rangle ,
\]

and this completes the proof of lemma 3.

DEFINITION 7. --- Let u be a continuous and positive endomorphism of the hilbertian space E. Let

\[
\operatorname{Tr} (u) = \sup _ {p \in \mathcal {F}} \tau (p u p)\tag{24}
\]

(upper bound in \(\overline{\mathbf{R}}_+\) ), where \(\mathcal{F}\) is the set of all finite rank orthoprojectors on E. We say that \(\operatorname{Tr}(u)\) is the trace of \(u\) .

Let \(p\) be the orthoprojector from \(E\) onto a finite dimensional vector subspace of \(E\) , and let \((x_1, \ldots, x_m)\) be an orthonormal basis of \(F\) . We have established the relation \(\tau(pup) = \sum_{i=1}^{m} \langle x_i | u(x_i) \rangle\) . Consequently, we can define the trace by the formula

\[
\operatorname{Tr} (u) = \sup _ {x _ {1}, \dots , x _ {m}} \sum_ {i = 1} ^ {m} \left\langle x _ {i} | u (x _ {i}) \right\rangle ,\tag{24'}
\]

where \((x_{1},\dots,x_{m})\) ranges over the set of all finite orthonormal sequences of vectors of E.

By lemma 3 (V, p.~48), we have

\[
\operatorname{Tr} (u) = \sum_ {i \in \mathbf {I}} \left\langle e _ {i} | u (e _ {i}) \right\rangle\tag{25}
\]

for every orthonormal basis \((e_{i})_{i\in I}\) of E. From this, we deduce

\[
\operatorname{Tr} (u + v) = \operatorname{Tr} (u) + \operatorname{Tr} (v)\tag{26}
\]

\[
\operatorname{Tr} (\lambda u) = \lambda . \operatorname{Tr} (u)\tag{27}
\]

for all continuous and positive endomorphisms u and v of E and for every real number \(\lambda \geqslant 0\) (we make the convention \(0. (+\infty) = 0\) in (27)). Let \(\phi\) be an isomorphism from E onto a hilbertian space F; since \(\phi\) transforms every orthonormal basis of E into an orthonormal basis of F, we get from (25) that

\[
\operatorname{Tr} \left(\phi u \phi^ {- 1}\right) = \operatorname{Tr} (u).\tag{28}
\]

Let \((u_{\alpha})_{\alpha \in \mathbf{A}}\) be a non-empty directed increasing and bounded family of continuous and positive endomorphisms of E; let \(u = \sup_{\alpha} u_{\alpha}\) , then \(\langle x|u(x) \rangle = \sup_{\alpha} \langle x|u_{\alpha}(x) \rangle\) for all \(x \in \mathrm{E}\) (V, p.~46, prop. 13). We have \(\operatorname{Tr}(u) = \sup_{\mathbf{J} \subset \mathbf{I}} \sum_{i \in \mathbf{J}} \langle e_i | u(e_i) \rangle\) , where \(\mathbf{J}\) ranges over all finite subsets of I, hence

\[
\operatorname{Tr} (u) = \sup _ {\alpha} \operatorname{Tr} (u _ {\alpha}) \quad \text { for } \quad u = \sup _ {\alpha} u _ {\alpha}.\tag{29}
\]

Let \(p_{\mathrm{F}}\) be the orthoprojector from E onto the hilbertian subspace F; there exists an orthonormal basis \((e_i)_{i \in \mathrm{I}}\) of E and a subset J of I, such that \((e_i)_{i \in \mathrm{J}}\) is an orthonormal basis of F. We have \(\operatorname{Tr}(p_{\mathrm{F}} up_{\mathrm{F}}) = \sum_{i \in \mathrm{J}} \langle e_i | u(e_i) \rangle\) . This formula has two consequences: firstly, we have \(\operatorname{Tr}(p_{\mathrm{F}} up_{\mathrm{F}}) \leqslant \operatorname{Tr}(u)\) ; secondly, taking \(u = 1_{\mathrm{E}}\) , we get

\[
\operatorname{Tr} (p _ {\mathrm{F}}) = \left\{ \begin{array}{l l} \dim \mathrm{F} & \text { if   F   is   finite   dimensional } \\ + \infty & \text { if   not. } \end{array} \right.\tag{30}
\]

DEFINITION 8. --- Let E be a complex hilbertian space. We write \(\mathcal{L}^{1}(\mathrm{E})\) for the vector subspace of \(\mathcal{L}(\mathrm{E})\) generated by all continuous, positive endomorphisms of E with finite trace.

By formula (25) of V, p.~50, the trace extends to a linear form on \(\mathcal{L}^1 (\mathrm{E})\) , again denoted by Tr, and satisfying the relation \(\operatorname{Tr}(u) = \sum_{i\in I}\langle e_i|u(e_i)\rangle\) for all \(u\) in \(\mathcal{L}^1 (\mathrm{E})\) and for every orthonormal basis \((e_i)_{i\in I}\) of E. For every \(u\in \mathcal{L}^1 (\mathrm{E})\) , we have \(u^{*}\in \mathcal{L}^{1}(\mathrm{E})\) and \(\operatorname {Tr}(u^{*}) = \overline{\operatorname{Tr}(u)}\) . Formula (28) of V, p.~50 extends to the case where \(u\) belongs to \(\mathcal{L}^1 (\mathrm{E})\) . Let F be a hilbertian subspace of E; by formula (30), the orthoprojector \(p_{\mathrm{F}}\) belongs to \(\mathcal{L}^1 (\mathrm{E})\) if and only if F is finite dimensional. For every \(a\) and \(b\) in E, we have \(4ab^{*} = \sum_{\varepsilon^{4} = 1}\varepsilon (a + \varepsilon b)(a + \varepsilon b)^{*}\) and \(cc^{*}\) is a positive operator with finite trace for all \(c\in \mathrm{E}\) : consequently, if \(u\) is a finite rank, continuous endomorphism of E, then \(u \in \mathcal{L}^1(\mathrm{E})\) and \(\operatorname{Tr}(u) = \tau(u)\) .

Let \(\mathbf{E}\) be a real hilbertian space, and let \(\mathrm{E}_{(\mathbf{C})}\) be its complexification (V, p.~5). We identify \(\mathbf{E}\) with a subset of \(\mathrm{E}_{(\mathbf{C})}\) . Then \(\mathcal{L}(\mathrm{E})\) can be identified with a real vector subspace of \(\mathcal{L}(\mathrm{E}_{(\mathbf{C})})\) consisting of all continuous linear mappings \(u\) from \(\mathrm{E}_{(\mathbf{C})}\) into \(\mathrm{E}_{(\mathbf{C})}\) such that \(u(\mathrm{E}) \subset \mathrm{E}\) . In this case we write \(\mathcal{L}^1(\mathrm{E}) = \mathcal{L}(\mathrm{E}) \cap \mathcal{L}^1(\mathrm{E}_{(\mathbf{C})})\) . For every \(u \in \mathcal{L}^1(\mathrm{E})\) , the trace \(\operatorname{Tr}(u)\) is real and is equal to \(\operatorname{Tr}(u^*)\) . Formulas (25) and (28) are again valid, \(\mathcal{L}_f(\mathrm{E}) \subset \mathcal{L}^1(\mathrm{E})\) and \(\operatorname{Tr}(u) = \tau(u)\) for all \(u \in \mathcal{L}_f(\mathrm{E})\) . Finally, a closed vector subspace \(\mathbf{F}\) of \(\mathbf{E}\) is finite dimensional if and only if \(p_{\mathrm{F}}\) belongs to \(\mathcal{L}^1(\mathrm{E})\) .

* Remark 1. --- We shall later define the notion of a nuclear mapping from a Banach space E into a Banach space F. We shall show that when \(\mathcal{L}^1 (\mathrm{E})\) consists of all nuclear mappings from E into E, then E is a real or complex hilbertian space. *

PROPOSITION 14. --- Let \(E_{1}, \ldots, E_{n}\) be hilbertian spaces, \(E = E_{1} \hat{\otimes}_{2} \ldots \hat{\otimes}_{2} E_{n}\) , and \(u_{i}\) a continuous endomorphism of \(E_{i}\) for \(1 \leqslant i \leqslant n\) . If \(u_{1}, \ldots, u_{n}\) are positive, then so is \(u = u_{1} \hat{\otimes}_{2} \ldots \hat{\otimes}_{2} u_{n}\) , and

\[
\operatorname{Tr} (u) = \prod_ {i = 1} ^ {n} \operatorname{Tr} (u _ {i}).\tag{31}
\]

If \(u_{i} \in \mathcal{L}^{1}(\mathrm{E}_{i})\) for all \(1 \leqslant i \leqslant n\) , then \(u \in \mathcal{L}^{1}(\mathrm{E})\) and formula (31) is again valid in this case.

Proceeding by induction on \(n\) , we immediately reduce to the case \(n = 2\) .

For i=1,2, we define a sesquilinear form \(\Phi_{i}\) on \(E_{i}\) by the formula \(\Phi_{i}(x,y)=\langle x|u_{i}(y)\rangle\) for x, y in \(E_{i}\) . If \(u_{1}\) and \(u_{2}\) are positive, the forms \(\Phi_{1}\) and \(\Phi_{2}\) are hermitian and positive. By prop. 1 of V, p.~25 there exists a positive hermitian form \(\Phi\) on the vector space \(E_{1}\otimes E_{2}\) such that

\[
\Phi (x _ {1} \otimes x _ {2}, y _ {1} \otimes y _ {2}) = \Phi_ {1} (x _ {1}, y _ {1}). \Phi_ {2} (x _ {2}, y _ {2})
\]

for \(x_{1}, y_{1}\) in \(E_{1}\) and \(x_{2}, y_{2}\) in \(E_{2}\) . We verify immediately the relation \(\Phi(z, t) = \langle z|u(t) \rangle\) for \(z\) and \(t\) in \(E_{1} \otimes E_{2}\) . Since \(\Phi\) is positive, we have \(\langle z|u(z) \rangle \geqslant 0\) for all \(z\) in \(E_{1} \otimes E_{2}\) . Since \(u\) is continuous and \(E_{1} \otimes E_{2}\) is dense in the hilbertian space \(E = E_{1} \hat{\otimes}_{2} E_{2}\) , we conclude that \(u\) is a continuous and positive endomorphism of \(E\) .

Let \((e_i)_{i\in I}\) be an orthonormal basis of \(\mathbf{E}_1\) and \((f_j)_{j\in J}\) an orthonormal basis of \(\mathbf{E}_2\) ; then the family \((e_i\otimes f_i)_{i\in I,j\in J}\) is an orthonormal basis of E and we have

\[
\begin{array}{r l} & {\mathrm{Tr} (u) = \sum_ {i \in \mathbf {I}} \sum_ {j \in \mathbf {J}} \langle e _ {i} \otimes f _ {j} | u (e _ {i} \otimes f _ {j}) \rangle} \\ & {\quad = \sum_ {i \in \mathbf {I}} \sum_ {j \in \mathbf {J}} \langle e _ {i} | u _ {1} (e _ {i}) \rangle . \langle f _ {j} | u _ {2} (f _ {j}) \rangle} \\ & {\quad = \mathrm{Tr} (u _ {1}). \mathrm{Tr} (u _ {2}).} \end{array}
\]

In particular, if \(u_{1}\) and \(u_{2}\) are positive endomorphisms with finite trace, then so is u. By linearity, we deduce that u belongs to \(\mathcal{L}^{1}(\mathrm{E})\) when K = C and that the \(u_{i}\) belong to \(\mathcal{L}^{1}(\mathrm{E}_{i})\) for i = 1, 2; formula (31) extends to this case by linearity. Finally, the case when K = R and the \(u_{i} \in \mathcal{L}^{1}(\mathrm{E}_{i})\) reduces to the complex case by extension of the scalars.

Remark 2. --- Let E be a hilbertian space, which is the hilbertian sum of a family \((\mathrm{E}_{i})_{i\in\mathrm{I}}\) of hilbertian subspaces. Let u be an element of \(\mathcal{L}(\mathrm{E})\) such that \(u(\mathrm{E}_{i}) \subset \mathrm{E}_{i}\) for all \(i \in I\) ; let \(u_{i}\) be the element of \(\mathcal{L}(\mathrm{E}_{i})\) which coincides with u on \(E_{i}\) . Then \(\operatorname{Tr}(u) = \sum_{i \in \mathrm{I}} \operatorname{Tr}(u_{i})\) when u is positive, or belongs to \(\mathcal{L}^{1}(\mathrm{E})\) ; this relation follows from formula (25) of V, p.~50 applied to an orthonormal basis of E which is the union of orthonormal bases of each of the \(E_{i}\) .

\subsection{Hilbert-Schmidt mappings}

DEFINITION 9. --- Let E and F be two hilbertian spaces. A continuous linear mapping u from E into F is called a Hilbert-Schmidt mapping if the trace of the positive endomorphism \(u^{*}u\) of E is finite. The set of all Hilbert-Schmidt mappings from E into F is denoted by \(\mathcal{L}^{2}(\mathrm{E},\mathrm{F})\) .

When \(\mathrm{E} = \mathrm{F}\) , we write \(\mathcal{L}^2(\mathrm{E})\) for \(\mathcal{L}^2(\mathrm{E};\mathrm{E})\) .

For every \(u \in \mathcal{L}(\mathrm{E},\mathrm{F})\) , let \(\| u \|_2 = \operatorname{Tr}(u^* u)^{1/2}\) , so that \(u\) belongs to \(\mathcal{L}^2(\mathrm{E};\mathrm{F})\) if and only if \(\| u \|_2\) is finite. By the definition of the trace, we get

\[
\| u \| _ {2} ^ {2} = \sup _ {x _ {1}, \dots , x _ {m}} \sum_ {i = 1} ^ {m} \left\| u (x _ {i}) \right\| ^ {2}\tag{32}
\]

where \((x_{1},\ldots ,x_{m})\) range over the set of finite orthonormal sequences in E. In particular, taking \(m = 1\) in formula (32), we have

\[
\| u \| \leqslant \| u \| _ {2} \quad (u \in \mathscr {L} (\mathrm{E}; \mathrm{F})).\tag{33}
\]

Let \((e_i)_{i\in I}\) be an orthonormal basis of \(E\) and \((f_j)_{j\in J}\) an orthonormal basis of \(F\) . By formula (25) of \(V\) , p.~50 and the Parseval's relation (V, p.~22), we have

\[
\left\| u \right\| _ {2} ^ {2} = \sum_ {i \in I} \left\| u (e _ {i}) \right\| ^ {2} = \sum_ {i, j} \left| \langle f _ {j} | u (e _ {i}) \rangle \right| ^ {2}.\tag{34}
\]

Since \(\left|\langle f_j | u(e_i) \rangle\right| = \left|\langle e_i | u^*(f_j) \rangle\right|\) , formula (34) implies that

\[
\| u ^ {*} \| _ {2} = \| u \|;\tag{35}
\]

hence, the adjoint of a Hilbert-Schmidt mapping is a Hilbert-Schmidt mapping. Let \(E_{1}\) , \(F_{1}\) be hilbertian spaces and \(v: E_{1} \to E\) , \(w: F \to F_{1}\) continuous linear mappings. From (32), we deduce immediately that

\[
\| w u \| _ {2} \leqslant \| w \|. \| u \| _ {2}.\tag{36}
\]

By (35), (36) and the relation \(uv = (v^{*}u^{*})^{*}\) , we get

\[
\left\| u v \right\| _ {2} \leqslant \left\| u \right\| _ {2} \left\| v \right\|.\tag{37}
\]

In particular, if \(u\) belongs to \(\mathcal{L}^2(\mathrm{E},\mathrm{F})\) then wuv belongs to \(\mathcal{L}^2(\mathrm{E}_1,\mathrm{F}_1)\) .

THEOREM 1. --- Let E and F be two hilbertian spaces.

\begin{enumerate}
\def\labelenumi{(\roman{enumi})}
\item
  The set \(\mathcal{L}^2 (\mathrm{E},\mathrm{F})\) is a vector subspace of \(\mathcal{L}(\mathrm{E};\mathrm{F})\) and \(u\mapsto \| u\| _2\) is a hilbertian norm (V, p.~6) on \(\mathcal{L}^2 (\mathrm{E};\mathrm{F})\) .
\item
  The isomorphism \(\theta\) from \(\mathbf{F} \otimes \mathbf{E}'\) onto \(\mathcal{L}_f(\mathbf{E};\mathbf{F})\) characterized by \(\theta(y \otimes x^*) = yx^*\) extends to an isomorphism \(\hat{\theta}\) from \(\mathbf{F} \hat{\otimes}_2 \mathbf{E}'\) onto \(\mathcal{L}^2(\mathbf{E};\mathbf{F})\) . In particular, \(\mathcal{L}_f(\mathbf{E};\mathbf{F})\) is dense in \(\mathcal{L}^2(\mathbf{E};\mathbf{F})\) .
\end{enumerate}

Let \((e_i)_{i\in \mathbf{I}}(\mathrm{resp.}(f_j)_{j\in \mathbf{J})}\) be an orthonormal basis of E (resp. F). For every \(u\in \mathcal{L}(\mathbf{E};\mathbf{F})\) let \(\Lambda (u)\) be the matrix of \(u\) with respect to chosen orthonormal bases for E and F (V, p.~22). Let \(\| \mathbf{a}\| _2\) denote the norm of an element \(a\) of the hilbertian space \(\ell^2 (\mathbf{J}\times \mathbf{I})\) . By formula (34), \(\Lambda\) is a mapping from \(\mathcal{L}^2 (\mathbf{E};\mathbf{F})\) into \(\ell^2 (\mathbf{J}\times \mathbf{I})\) such that \(\| \Lambda (u)\| _2 = \| u\|\) ; it is clear that \(\Lambda\) is injective. To prove (i), it is enough to prove that \(\Lambda\) is surjective. Let \(\mathbf{a} = (a_{ji})\) be an element of \(\ell^2 (\mathbf{J}\times \mathbf{I})\) ; by Cauchy-Schwarz inequality, we have

\[
\left| \sum_ {j, i} \overline {{{{\eta}}}} _ {j} a _ {j i} \xi_ {i} \right| ^ {2} \leqslant \sum_ {j, i} | a _ {j i} | ^ {2} \sum_ {j, i} | \overline {{{{\eta}}}} _ {j} \xi_ {i} | ^ {2} = \| \mathbf {a} \| _ {2} ^ {2} \| \boldsymbol {\xi} \| ^ {2} \| \boldsymbol {\eta} \| ^ {2}
\]

for every \(\xi = (\xi_i)\) in \(\ell^2 (\mathrm{I})\) and \(\pmb{\eta} = (\eta_j)\) in \(\ell^2 (\mathrm{J})\) . Then there exists a continuous sesquilinear form \(\Phi\) on \(\mathbf{F}\times \mathbf{E}\) such that \(\Phi (y,x) = \sum_{j,i}\overline{\eta}_ja_{ji}\xi_i\) for \(x = \sum_{i}\xi_{i}e_{i}\) in E and \(y = \sum_{j}\eta_{j}f_{j}\) in F. Let \(u\in \mathcal{L}(\mathrm{E};\mathrm{F})\) be such that \(\Phi (y,x) = \langle y|u(x)\rangle\) (V, p.~16, cor. 2). We get

\[
a _ {j i} = \Phi (f _ {j}, e _ {i}) = \langle f _ {j} | u (e _ {i}) \rangle \quad \text { for } \quad i \in \mathrm{I}, j \in \mathrm{J}  ,
\]

hence \(\mathbf{a} = \Lambda(u)\) .

Since \(\Lambda\) is a hilbertian space isomorphism from \(\mathcal{L}^2 (\mathrm{E};\mathrm{F})\) onto \(\ell^2 (\mathrm{J}\times \mathrm{I})\) and since \((f_j\otimes e_i^*)\) is an orthonormal basis of \(\mathbf{F}\hat{\otimes}_2\mathbf{E}'\) , there exists an isomorphism \(\hat{\theta}\) from \(\mathbf{F}\hat{\otimes}_2\mathbf{E}'\) onto \(\mathcal{L}^2 (\mathrm{E};\mathrm{F})\) such that

\[
\langle f _ {j} | \hat {\theta} (t) e _ {i} \rangle = \langle f _ {j} \otimes e _ {i} ^ {*} | t \rangle
\]

for every \(i \in \mathbf{I}\) , \(j \in \mathbf{I}\) and \(t \in \mathbf{F} \hat{\otimes}_2 \mathbf{E}'\) . In particular, for \(t = y \otimes x^*\) , we find

\[
\langle f _ {j} | \hat {\theta} (y \otimes x ^ {*}) e _ {i} \rangle = \langle f _ {j} \otimes e _ {i} ^ {*} | y \otimes x ^ {*} \rangle = \langle f _ {j} | y \rangle \langle x | e _ {i} \rangle = \langle f _ {j} | y x ^ {*} e _ {i} \rangle
\]

hence \(\hat{\theta}(y\otimes x^{*}) = yx^{*}\) . This proves (ii).

Q.E.D.

Examples. --- 1) Let I and J be two sets. By the proof given above, in order that a mapping \(u\) from \(\ell^2(\mathbf{I})\) into \(\ell^2(\mathbf{J})\) be a Hilbert-Schmidt mapping, it is necessary and sufficient that there exists a matrix \((a_{ji})\) in \(\ell^2(\mathbf{J} \times \mathbf{I})\) such that \(u(\xi)_j = \sum_{i \in \mathbf{I}} a_{ji} \xi_i\) for all \(\xi = (\xi_i)\) in \(\ell^2(\mathbf{I})\) .

* 2) Let X and Y be two Hausdorff topological spaces, endowed respectively with positive measures \(\mu\) and \(\nu\) . We can show that the Hilbert-Schmidt mappings from \(\mathcal{L}^2(\mathrm{X})\) into \(\mathcal{L}^2(\mathrm{Y})\) correspond bijectively to classes of square integrable functions on \(\mathrm{Y} \times \mathrm{X}\) ; to the class of a function \(\mathrm{N} \in \mathcal{L}^2(\mathrm{Y} \times \mathrm{X}, \nu \otimes \mu)\) corresponds the mapping \(u_{\mathrm{N}}\) given by

\[
(u _ {\mathrm{N}} f) (y) = \int_ {\mathrm{X}} \mathrm{N} (y, x) f (x) d \mu (x)\tag{38}
\]

for v-almost all \(y \in \mathbf{Y}\) and \(f \in \mathcal{L}^2(\mathbf{X}, \mu)\) . We have

\[
\| u _ {\mathrm{N}} \| _ {2} ^ {2} = \int_ {\mathrm{X}} \int_ {\mathrm{Y}} | \mathrm{N} (y, x) | ^ {2} d \mu (x) d v (y). *\tag{39}
\]

Remarks. --- 1) Suppose \(\mathbf{K} = \mathbf{C}\) . Let \(u\) and \(v\) be in \(\mathcal{L}^2(\mathrm{E};\mathrm{F})\) . We have the relation \(4u^{*}v = \sum_{\varepsilon^{4} = 1}\overline{\varepsilon} (u + \varepsilon v)^{*}(u + \varepsilon v)\) , hence \(u^{*}v\) belongs to \(\mathcal{L}^1(\mathrm{E})\) . The scalar product in the hilbertian space \(\mathcal{L}^2(\mathrm{E};\mathrm{F})\) is given by

\[
\langle u | v \rangle = \operatorname{Tr} (u ^ {*} v)\tag{40}
\]

since this formula defines a hermitian form on \(\mathcal{L}^{2}(\mathrm{E};\mathrm{F})\) and we get \(\langle u|u\rangle=\|u\|_{2}^{2}\) . If \(u\in\mathcal{L}^{2}(\mathrm{E};\mathrm{F})\) and \(v\in\mathcal{L}^{2}(\mathrm{F};\mathrm{E})\) , then vu belongs to \(\mathcal{L}^{1}(\mathrm{E})\) and uv to \(\mathcal{L}^{1}(\mathrm{F})\) by the preceding; moreover, we have

\[
\operatorname{Tr} (u v) = \operatorname{Tr} (v u).\tag{41}
\]

By linearity and continuity, it is enough to verify this formula when \(u = y_{1}x_{1}^{*}\) and \(v = x_{2}y_{2}^{*}\) (with \(x_{1}, x_{2}\) in E, \(y_{1}, y_{2}\) in F); but then uv is the mapping \(y \mapsto y_{1} \langle x_{1}|x_{2} \rangle \langle y_{2}|y \rangle\) and vu the mapping \(x \mapsto x_{2} \langle y_{2}|y_{1} \rangle \langle x_{1}|x \rangle\) , and (41) follows from formula (22) of V, p.~48.

Consequently, if \(u_{1}, u_{2}\) are two elements of \(\mathcal{L}^2(\mathrm{E};\mathrm{F})\) , we have, in the hilbertian space \(\mathcal{L}^2(\mathrm{F};\mathrm{E})\) ,

\[
\langle u _ {1} ^ {*} | u _ {2} ^ {*} \rangle = \operatorname{Tr} (u _ {1} u _ {2} ^ {*}) = \operatorname{Tr} (u _ {2} ^ {*} u _ {1}) = \langle u _ {2} | u _ {1} \rangle = \overline {{\langle u _ {1} | u _ {2} \rangle}};\tag{42}
\]

in other words, \(u \mapsto u^*\) is an isomorphism from the hilbertian space \(\mathcal{L}^2(\mathrm{E};\mathrm{F})\) onto the conjugate (V, p.~6) of the hilbertian space \(\mathcal{L}^2(\mathrm{F};\mathrm{E})\) . If we identify this conjugate with the dual of \(\mathcal{L}^2 (\mathrm{F};\mathrm{E})\) (V, p.~15), we see that \(\mathcal{L}^2 (\mathrm{E};\mathrm{F})\) can be identified with the dual of \(\mathcal{L}^2 (\mathrm{F};\mathrm{E})\) , the canonical bilinear form \((v,u)\mapsto \langle v,u\rangle\) being identified with \((v,u)\mapsto \operatorname {Tr}(vu)\) .

\begin{enumerate}
\def\labelenumi{\arabic{enumi})}
\setcounter{enumi}{1}
\tightlist
\item
  Suppose K = R. We leave it to the reader to verify that formulas (40) and (41) are again valid, and to show that \(\mathcal{L}^{2}(E; F)\) can be identified with the dual of \(\mathcal{L}^{2}(F; E)\) by means of the bilinear form \((u, v) \mapsto \operatorname{Tr}(uv)\) .
\end{enumerate}

\subsection{Diagonalization of Hilbert-Schmidt mappings}

THEOREM 2. --- Let E and F be two hilbertian spaces and u a Hilbert-Schmidt mapping from E into F. There exists an orthonormal basis \((e_{i})_{i\in I}\) of E which is transformed by u into an orthogonal family in F.

Let B denote the (closed) unit ball of E, with the weakened topology assigned to it; this is a compact space (V, p.~17). We put \(\mathbf{Q}(x) = \|u(x)\|^{2}\) for all \(x \in B\) . Finally let P denote the set of all vectors x in E satisfying the following property;

\begin{enumerate}
\def\labelenumi{(\Alph{enumi})}
\setcounter{enumi}{7}
\tightlist
\item
  For every \(y \in \mathbf{E}\) orthogonal to \(x\) , the element \(u(y)\) of \(\mathbf{E}\) is orthogonal to \(u(x)\) .
\end{enumerate}

Lemma 4. --- The function \(\mathbf{Q}:\mathbf{B}\to \mathbf{R}\) is continuous.

Let \((f_j)_{j\in \mathbf{J}}\) be an orthonormal basis of F. Put \(\lambda_{j} = \| u^{*}(f_{j})\|^{2}\) for all \(j\in \mathbf{J}\) . Since \(u\) belongs to \(\mathcal{L}^2 (\mathrm{E};\mathrm{F})\) we have \(u^{*}\in \mathcal{L}^{2}(\mathrm{F};\mathrm{E})\) , hence \(\sum_{j}\lambda_{j} < + \infty\) . Further, we have

\[
Q (x) = \left\| u (x) \right\| ^ {2} = \sum_ {j} \left| \langle u ^ {*} (f _ {j}) | x \rangle \right| ^ {2}\tag{43}
\]

by Parseval's formula (V, p.~22) and the definition of the adjoint (V, p.~38). For every \(x \in \mathbf{B}\) , \(|\langle u^{*}(f_{j})|x\rangle|^{2} \leqslant \lambda_{j}\) by Cauchy-Schwarz inequality; consequently, the convergence of the sum in formula (43) is uniform on B, hence lemma 4 (GT, X, § 1, No.~6).

Lemma 5. --- Let \(\mathbf{E}_1\) be a closed vector subspace of \(\mathbf{E}\) , stable under \(u^*u\) . If \(\mathbf{E}_1 \neq \{0\}\) , then there exists a vector of norm 1 in \(\mathbf{E}_1 \cap \mathbf{P}\) .

Since B is weakly compact, so is the weakly closed subspace \(B \cap E_{1}\) of B. Hence there exists (GT, IV, §6, No.~1, th. 1) a point \(x_{0}\) in \(B \cap E_{1}\) such that \(Q(x_{0} \geqslant Q(x))\) for all \(x \in B \cap E_{1}\) . If \(Q(x_{0}) = 0\) , we have \(Q(x) = 0\) and so \(u(x) = 0\) for all \(x \in B \cap E_{1}\) . Thus \(E_{1} \subset P\) and lemma 5 follows in this case.

Suppose now that \(\mathbf{Q}(x_0) > 0\) , then \(x_0 \neq 0\) . Since the vector \(\| x_0 \|^{-1} \cdot x_0\) belongs to \(\mathbf{B} \cap \mathbf{E}_1\) , we have

\[
\mathrm{Q} \left(x _ {0}\right) \geqslant \mathrm{Q} \left(\left\| x _ {0} \right\| ^ {- 1}. x _ {0}\right) = \mathrm{Q} \left(x _ {0}\right) / \left\| x _ {0} \right\| ^ {2}
\]

i.e.~\(\|x_{0}\|=1\) . We shall prove that \(x_{0}\) belongs to P; let \(y\in E\) be orthogonal to \(x_{0}\) . It is enough to prove that \(u(y)\) is orthogonal to \(u(x_{0})\) . But since y is the sum of a vector of \(E_{1}\) and a vector orthogonal to \(E_{1}\) , and both orthogonal to \(x_{0}\) (since \(x_{0}\in E_{1}\) ), it is enough to consider the following two cases:

\begin{enumerate}
\def\labelenumi{\alph{enumi})}
\item
  \(y\) is orthogonal to \(\mathbf{E}_1\) : since \(\mathbf{E}_1\) is stable under \(u^* u, u^* u(x_0) \in \mathbf{E}_1\) , hence \(0 = \langle y | u^* u(x_0) \rangle = \langle u(y) | u(x_0) \rangle\) .
\item
  \(y\) belongs to \(\mathbf{E}_1\) : for all \(t \in \mathbf{R}\) , the vector \(x(t) = (x_0 + ty) / \| x_0 + ty \|\) belongs to \(\mathbf{B} \cap \mathbf{E}_1\) . We have \(Q(xt) = f(t) / g(t)\) with
\end{enumerate}

\[
f (t) = \left\| u (x _ {0}) \right\| ^ {2} + 2 t \mathscr {R} \left\langle u (x _ {0}) | u (y) \right\rangle + t ^ {2} \left\| u (y) \right\| ^ {2}
\]

\[
g (t) = 1 + t ^ {2} \| y \| ^ {2}.
\]

In view of the definition of \(x_0\) , we have \(\mathrm{Q}(x(0)) \geqslant \mathrm{Q}(x(t))\) for all real \(t\) , hence \(\frac{d}{dt} \mathrm{Q}(x(t))\) is zero for \(t = 0\) . But \(f(0) = \| u(x_0)\|^2\) , \(g(0) = 1\) , \(f'(0) = 2\mathcal{R}\langle u(x_0)|u(y)\rangle\) , \(g'(0) = 0\) . Since

\[
\frac {d}{d t} \mathrm{Q} (x (t)) = \frac {f ^ {\prime} (t) g (t) - f (t) g ^ {\prime} (t)}{g (t) ^ {2}},
\]

we conclude that \(f'(0)=0\) , that is, \(\mathcal{R}\langle u(x_{0})|u(y)\rangle=0\) . When K=R, \(u(x_{0})\) is orthogonal to \(u(y)\) , when K=C, the vector iy belongs to \(E_{1}\) and is orthogonal to \(x_{0}\) , hence \(\mathcal{I}\langle u(x_{0})|u(y)\rangle=-\mathcal{R}\langle u(x_{0})|u(iy)\rangle=0\) , and finally \(u(x_{0})\) is orthogonal to \(u(y)\) . This proves lemma 5.

Now we prove th. 2. Applying th. 1 of S, III, §4, No.~5 we see, as in V, p.~23, that there exists a set S which is maximal among the orthonormal subsets of E contained in P. Let \(E_{1}\) be the set of all vectors orthogonal to S. Let \(y \in E_{1}\) ; if \(x \in S\) , the vectors x and y are orthogonal, and since \(S \subset P\) , we conclude that \(u(x)\) and \(u(y)\) are orthogonal; then

\[
\langle x | u ^ {*} u (y) \rangle = \langle u (x) | u (y) \rangle = 0
\]

and \(u^{*}u(y)\) is orthogonal to S. Hence \(E_{1}\) is stable under \(u^{*}u\) . If we had \(E_{1} \neq \{0\}\) , there would exist a vector \(x\) of norm 1 in \(E_{1} \cap P\) (lemma 5) and \(S \cup \{x\}\) would be an orthonormal subset of E contained in P. This would contradict the maximal character of S. Hence \(E_{1} = \{0\}\) and S is an orthonormal basis of E. Q.E.D.

COROLLARY 1. --- Let v be a continuous, positive endomorphism with finite trace of the hilbertian space E. There exists an orthonormal basis \((e_{i})_{i\in I}\) of E and a summable family of positive real numbers \((\lambda_{i})_{i\in I}\) such that \(v(e_{i}) = \lambda_{i}e_{i}\) for all \(i \in I\) .

Let \(\Phi(x, y) = \langle x | v(y) \rangle\) for \(x, y\) in E. Then \(\Phi\) is a positive hermitian form on E. There exists (V, p.~8, corollary) a hilbertian space F and a continuous linear mapping \(u\) from E into F such that \(\Phi(x, y) = \langle u(x) | u(y) \rangle\) for \(x, y\) in E. In other words, we have \(v = u^{*}u\) . By virtue of def. 9 (V, p.~52), \(u\) is a Hilbert-Schmidt mapping from E into F. By th. 2, there exists an orthonormal basis \((e_{i})_{i \in I}\) of E such that the vectors \(u(e_{i})\) are two by two orthogonal. Let \(i \in I\) ; for every \(j \neq i\) in I, we have

\[
\langle e _ {j} | v (e _ {i}) \rangle = \langle u (e _ {j}) | u (e _ {i}) \rangle = 0
\]

hence \(v(e_{i})\) is proportional to \(e_{i}\) and is of the form \(\lambda_{i}e_{i}\) , where \(\lambda_{i}=\langle e_{i}|v(e_{i})\rangle\) ; then

\[
\lambda_ {i} \geqslant 0 \quad \text { and } \quad \sum_ {i \in I} \lambda_ {i} = \operatorname{Tr} (v) <   + \infty .
\]

COROLLARY 2. --- Let E be a hilbertian space. Then \(\mathcal{L}^{1}(\mathrm{E})\subset \mathcal{L}^{2}(\mathrm{E})\)

The real case reduces to the complex case by the extension of scalars; we can therefore assume that \(\mathbf{K} = \mathbf{C}\) .

Since \(\mathcal{L}^{2}(E)\) is a vector subspace of \(\mathcal{L}(E)\) , it is enough to prove that every continuous and positive endomorphism v of E with finite trace belongs to \(\mathcal{L}^{2}(E)\) . With the notations of cor. 1, we have

\[
\sum_ {i \in \mathbf {I}} \left\| v (e _ {i}) \right\| ^ {2} = \sum_ {i \in \mathbf {I}} \lambda_ {i} ^ {2} \leqslant (\sum_ {i} \lambda_ {i}) ^ {2} <   + \infty .
\]

COROLLARY 3. --- Let v be a continuous positive endomorphism of the hilbertian space E with a finite trace. There exists a positive Hilbert-Schmidt endomorphism w of E such that \(v = w^{2}\) and such that v commutes with w.

With the notations of cor. 1, it is enough to consider the endomorphism \(w\) which transforms the vector \(\sum_{i\in I}\xi_i e_i\) into the vector \(\sum_{i}\lambda_i^{1 / 2}\xi_i e_i\) .

Remark. --- With the notations of th. 2, let J be the set of all \(i \in \mathbf{I}\) such that \(u(e_i) \neq 0\) . For all \(i \in \mathbf{J}\) , let \(\lambda_i = \| u(e_i)\|\) and \(f_i = \lambda_i^{-1}u(e_i)\) . Then \((e_i)_{i \in \mathbf{J}}(\text{resp. } (f_i)_{i \in \mathbf{J})}\) is an orthonormal basis of the initial (resp. final) subspace of \(u\) , we have \(u(e_i) = \lambda_i f_i\) for all \(i \in \mathbf{J}\) and \(\sum_{i \in \mathbf{J}} \lambda_i^2 = \| u\|_2^2\) is finite.

\subsection{Trace of a quadratic form with respect to another}

In this section, E will denote a real vector space and Q, H two positive quadratic forms on E. There exist two symmetric bilinear forms \((x, y) \mapsto \langle x | y \rangle_{\mathbf{Q}}\) and \((x, y) \mapsto \langle x | y \rangle_{\mathbf{H}}\) on \(E \times E\) , characterized by

\[
\mathrm{Q} (x) = \langle x | x \rangle_ {\mathrm{Q}}, \quad \mathrm{H} (x) = \langle x | x \rangle_ {\mathrm{H}}
\]

for all \(x \in \mathbf{E}\) .

We call the trace of Q with respect to H, and write \(\mathrm{Tr}(\mathrm{Q}/\mathrm{H})\) , a real positive number, finite or not, defined as follows :

\begin{enumerate}
\def\labelenumi{\alph{enumi})}
\tightlist
\item
  If there exists \(x \in \mathbf{E}\) with \(\mathbf{H}(x) = 0\) and \(\mathbf{Q}(x) \neq 0\) , we put \(\mathrm{Tr}(\mathbf{Q} / \mathbf{H}) = +\infty\) .
\item
  Otherwise, \(\mathrm{Tr}(\mathbf{Q} / \mathbf{H})\) is the upper bound of the set of all numbers of the form \(\sum_{i=1}^{m} \mathbf{Q}(x_i)\) where \((x_1, \ldots, x_m)\) range over the set of finite sequences of elements of \(\mathbf{E}\) such that \(\langle x_i | x_j \rangle_{\mathbf{H}} = \delta_{ij}\) (Kronecker's symbol).
\end{enumerate}

Remarks. --- 1) For every subspace F of E, let \(\mathrm{Q}_{\mathrm{F}}\) denote the restriction of Q to F and \(\mathrm{H}_{\mathrm{F}}\) that of H. We have \(\mathrm{Tr}(\mathrm{Q}_{\mathrm{F}} / \mathrm{H}_{\mathrm{F}}) \leqslant \mathrm{Tr}(\mathrm{Q} / \mathrm{H})\) and \(\mathrm{Tr}(\mathrm{Q} / \mathrm{H})\) is the upper bound of the set of all numbers \(\mathrm{Tr}(\mathrm{Q}_{\mathrm{F}} / \mathrm{H}_{\mathrm{F}})\) where F ranges over the family of all finite dimensional vector subspaces of E.

\begin{enumerate}
\def\labelenumi{\arabic{enumi})}
\setcounter{enumi}{1}
\tightlist
\item
  Let \(\mathbf{E}_1\) be a real vector space, \(\mathbf{Q}_1\) and \(\mathbf{H}_1\) two positive quadratic forms on \(\mathbf{E}_1\) and \(\pi: \mathbf{E} \to \mathbf{E}_1\) a surjective linear mapping. If \(\mathbf{Q} = \mathbf{Q}_1 \circ \pi\) and \(\mathbf{H} = \mathbf{H}_1 \circ \pi\) , then \(\operatorname{Tr}(\mathbf{Q}/\mathbf{H}) = \operatorname{Tr}(\mathbf{Q}_1/\mathbf{H}_1)\) .
\end{enumerate}

PROPOSITION 15. --- Suppose that there exists a real hilbertian space structure on E such that \(\mathbf{H}(x) = \|x\|^2\) for all \(x \in E\) . For \(\operatorname{Tr}(\mathbf{Q}/\mathbf{H})\) to be finite, it is necessary and sufficient that there exists a continuous and positive endomorphism u of E with finite trace, such that \(\mathbf{Q}(x) = \langle x|u(x) \rangle\) for all \(x \in E\) ; this endomorphism u is unique, and we have

\[
\operatorname{Tr} (u) = \operatorname{Tr} (\mathrm{Q} / \mathrm{H}) = \sum_ {i \in \mathrm{I}} \mathrm{Q} \left(e _ {i}\right)\tag{44}
\]

for every orthonormal bases \((e_i)_{i\in \mathbf{I}}\) of E.

Suppose that \(\operatorname{Tr}(\mathbf{Q} / \mathbf{H})\) is finite. For every \(x\in E\) of norm 1, we have \(\mathrm{H}(x) = 1\) , hence \(\mathrm{Q}(x)\leqslant \mathrm{Tr}(\mathrm{Q} / \mathrm{H})\) . Therefore, \(\mathrm{Q}(x)\leqslant \mathrm{Tr}(\mathrm{Q} / \mathrm{H}).\| x\|^2\) for all \(x\in E\) , and

\[
\left| \langle x | y \rangle_ {\mathrm{Q}} \right| \leqslant \mathrm{Q} (x) ^ {1 / 2} \mathrm{Q} (y) ^ {1 / 2} \leqslant \operatorname{Tr} (\mathrm{Q} / \mathrm{H}). \| x \|. \| y \|
\]

by the Cauchy-Schwarz inequality. Consequently, the bilinear form \((x,y)\mapsto\langle x|y\rangle_{\mathbf{Q}}\) on \(E\times E\) is continuous. There exists (V, p.~16, cor. 2) a mapping \(u\in\mathcal{L}(E)\) such that \(\langle x|y\rangle_{\mathbf{Q}}=\langle x|u(y)\rangle\) . We have \(\langle x|y\rangle_{\mathbf{Q}}=\langle y|x\rangle_{\mathbf{Q}}\) for x, y in E, hence u is hermitian; and \(\langle x|u(x)\rangle=\mathbf{Q}(x)\geqslant0\) , hence u is positive.

Conversely, let \(u\) be a continuous and positive endomorphism of \(E\) such that \(Q(x) = \langle x|u(x) \rangle\) for all \(x \in E\) . Then

\[
\langle x | u (y) \rangle = \frac {1}{2} (\mathrm{Q} (x + y) - \mathrm{Q} (x) - \mathrm{Q} (y)) = \langle x | y \rangle_ {\mathrm{Q}},
\]

which gives the uniqueness of u. By formula (24') (V, p.~50), we get

\[
\operatorname{Tr} (u) = \sup _ {x _ {1}, \dots , x _ {m}} \sum_ {i = 1} ^ {m} \left\langle x _ {i} \mid u \left(x _ {i}\right) \right\rangle = \sup _ {x _ {1}, \dots , x _ {m}} \sum_ {i = 1} ^ {m} Q \left(x _ {i}\right),
\]

where \((x_{1},\ldots ,x_{m})\) range over the set of all finite orthonormal sequences of elements of E. By the definition of \(\mathrm{Tr}(\mathrm{Q / H})\) , we get \(\mathrm{Tr}(u) = \mathrm{Tr}(\mathrm{Q / H})\) . Finally, for every orthonormal basis \((e_i)_{i\in \mathbf{I}}\) of E, we have \(\mathrm{Tr}(u) = \sum_{i\in \mathbf{I}}\langle e_i|u(e_i)\rangle\) by formula (25) of V, p.~50, hence \(\mathrm{Tr}(u) = \sum_{i\in \mathbf{I}}\mathrm{Q}(e_i)\) . Q.E.D.

Remarks. --- 3) Let E and F be two hilbertian spaces and v a linear, not necessarily continuous mapping from E into F. Let \(\mathbf{H}(x)=\|x\|^{2}\) and \(\mathbf{Q}(x)=\|v(x)\|^{2}\) for all \(x\in E\) . It follows from prop. 15 that v is a Hilbert-Schmidt mapping if and only if \(\operatorname{Tr}(\mathbf{Q}/\mathbf{H})\) is finite, and then \(\operatorname{Tr}(\mathbf{Q}/\mathbf{H})=\|v\|_{2}^{2}\) .

\begin{enumerate}
\def\labelenumi{\arabic{enumi})}
\setcounter{enumi}{3}
\tightlist
\item
  Suppose E is finite dimensional. When the quadratic form H is invertible, prop. 15 applies. Let \((e_{1}, \ldots, e_{n})\) be a basis of E. Put \(q_{ij} = \langle e_{i}|e_{j} \rangle_{Q}\) and \(h_{ij} = \langle e_{i}|e_{j} \rangle_{H}\) and introduce the matrices \(q = (q_{ij})\) and \(h = (h_{ij})\) . Let u be an endomorphism of E such that \(Q(x) = \langle x|u(x) \rangle_{H}\) for all \(x \in E\) . We have
\end{enumerate}

\[
\langle x | y \rangle_ {\mathrm{Q}} = \langle x | u (y) \rangle_ {\mathrm{H}} \quad (x, y \in \mathrm{E}),
\]

and hence the matrix of \(u\) with respect to the basis \((e_1, \ldots, e_n)\) of \(E\) is equal to \(h^{-1}q\) . By prop. 15, we have

\[
\operatorname{Tr} (\mathrm{Q} / \mathrm{H}) = \operatorname{Tr} (h ^ {- 1} q) = \operatorname{Tr} (q h ^ {- 1}).\tag{45}
\]

If the basis \((e_1, \dots, e_n)\) is orthonormal for \(\mathbf{H}\) , then \(h\) is the unit matrix of order \(n\) , and we get

\[
\operatorname{Tr} (\mathrm{Q} / \mathrm{H}) = \operatorname{Tr} (q) = \sum_ {i = 1} ^ {n} \mathrm{Q} (e _ {i});
\]

so that we get formula (44) in this case.

Now suppose that the quadratic form \(\mathbf{H}\) is not invertible. Let \(\mathbf{N}\) be the kernel of \(\mathbf{H}\) , and let \(\pi\) be the canonical mapping from \(\mathrm{E}\) onto \(\mathrm{E} / \mathrm{N}\) . There exists an invertible quadratic form \(\mathrm{H}_1\) on \(\mathrm{E} / \mathrm{N}\) such that \(\mathrm{H} = \mathrm{H}_1 \circ \pi\) . Let \((e_1, \dots, e_n)\) be a sequence of elements of \(\mathrm{E}\) such that the sequence \((\pi(e_1), \dots, \pi(e_m))\) is a basis of \(\mathrm{E} / \mathrm{N}\) , which is orthonormal for \(\mathrm{H}_1\) . Let \((e_1, \dots, e_n)\) be a basis of \(\mathbf{N}\) . Then \((e_1, \dots, e_n)\) is a basis of \(\mathrm{E}\) and we have

\[
\mathrm{H} \left(\xi_ {1} e _ {1} + \dots + \xi_ {n} e _ {n}\right) = \xi_ {1} ^ {2} + \dots + \xi_ {m} ^ {2}\tag{46}
\]

for all real numbers \(\xi_1, \ldots, \xi_n\) .

Suppose that for all \(x \in \mathbf{E}\) , the relation \(\mathrm{H}(x) = 0\) implies \(\mathrm{Q}(x) = 0\) ; in other words, suppose that there exists a quadratic form \(\mathbf{Q}_1\) on \(\mathrm{E} / \mathrm{N}\) such that \(\mathrm{Q} = \mathrm{Q}_1 \circ \pi\) . By remark 2 and prop. 15, we have,

\[
\operatorname{Tr} (\mathrm{Q} / \mathrm{H}) = \mathrm{Q} (e _ {1}) + \dots + \mathrm{Q} (e _ {m}).\tag{47}
\]

\exercisesheading
